\section{QWeb Templates}

\begin{frame}[fragile]{What Is QWeb?}
	\begin{itemize}
		\item \textbf{QWeb} is Odoo's XML templating engine.
		\item It renders HTML from XML templates + Python values.
		\item Used for:
		\begin{itemize}
			\item PDF/HTML reports
			\item website snippets and pages
			\item some portal and mail rendering flows
		\end{itemize}
\begin{block}{Core Idea}
	Template + values $\Rightarrow$ rendered HTML.
\end{block}
	\end{itemize}
\end{frame}

\begin{frame}[fragile]{Minimal QWeb Report Template}
\begin{lstlisting}[style=xml]
<template id="report_student_card">
  <t t-call="web.html_container">
	<t t-foreach="docs" t-as="o">
	  <t t-call="web.external_layout">
		<div class="page">
		  <h2>Student Card</h2>
		  <p><strong>Name:</strong> <t t-esc="o.name"/></p>
		  <p><strong>Code:</strong> <t t-esc="o.code"/></p>
		  <p><strong>Classroom:</strong>
			 <t t-esc="o.classroom_id.display_name"/>
		  </p>
		</div>
	  </t>
	</t>
  </t>
</template>
\end{lstlisting}
\end{frame}

\begin{frame}[fragile]{Breaking Down Key QWeb Directives}
	\begin{itemize}
		\item \textbf{\texttt{t-esc}}: escape and print a value (safe HTML escaping)
		\item \textbf{\texttt{t-out}}: modern output directive (preferred in newer code)
		\item \textbf{\texttt{t-raw}}: print raw HTML (use carefully)
		\item \textbf{\texttt{t-if}} / \textbf{\texttt{t-elif}} / \textbf{\texttt{t-else}}: conditions
		\item \textbf{\texttt{t-foreach}} / \textbf{\texttt{t-as}}: loops
		\item \textbf{\texttt{t-call}}: include/call another template
		\item \textbf{\texttt{t-set}} / \textbf{\texttt{t-value}}: create variables
	\end{itemize}
\end{frame}

\begin{frame}[fragile]{Conditions and Loops}
\begin{lstlisting}[style=xml]
<t t-if="o.email">
  <p>Email: <t t-esc="o.email"/></p>
</t>
<t t-else="">
  <p>No email on file</p>
</t>

<ul>
  <t t-foreach="o.attendance_ids" t-as="line">
	<li>
	  <t t-esc="line.date"/> -
	  <t t-esc="line.status"/>
	</li>
  </t>
</ul>
\end{lstlisting}
\end{frame}

\begin{frame}[fragile]{Variables with \texttt{t-set}}
\begin{lstlisting}[style=xml]
<t t-set="full_name" t-value="o.name"/>
<t t-set="is_adult" t-value="o.age &gt;= 18"/>

<p>Student: <t t-esc="full_name"/></p>
<p t-if="is_adult">This student is an adult.</p>
\end{lstlisting}
	\begin{itemize}
		\item Keep expressions simple and readable.
		\item Heavy business logic belongs in Python, not deep inside templates.
	\end{itemize}
\end{frame}

\begin{frame}[fragile]{Layout Helpers}
	\begin{itemize}
		\item Common report wrappers:
		\begin{itemize}
			\item \texttt{web.html\_container}
			\item \texttt{web.external\_layout}
			\item \texttt{web.internal\_layout}
		\end{itemize}
	\end{itemize}
\begin{lstlisting}[style=xml]
<t t-call="web.html_container">
  <t t-foreach="docs" t-as="o">
	<t t-call="web.external_layout">
	  <div class="page">
		<!-- report body -->
	  </div>
	</t>
  </t>
</t>
\end{lstlisting}
	\begin{itemize}
		\item \texttt{external\_layout} adds company header/footer.
	\end{itemize}
\end{frame}

\begin{frame}[fragile]{\texttt{docs} and Report Context}
	\begin{itemize}
		\item In reports, \texttt{docs} is the recordset being printed.
		\item Each record is usually looped as \texttt{o}.
\begin{lstlisting}[style=xml]
<t t-foreach="docs" t-as="o">
  <p><t t-esc="o.name"/></p>
</t>
\end{lstlisting}
		\item Extra values can come from a custom report parser / abstract model.
	\end{itemize}
\end{frame}

\begin{frame}[fragile]{Styling QWeb Reports}
\begin{lstlisting}[style=xml]
<div class="page">
  <h2 class="text-center">Student Report</h2>
  <table class="table table-sm o_main_table">
	<thead>
	  <tr>
		<th>Name</th>
		<th>Age</th>
		<th>Classroom</th>
	  </tr>
	</thead>
	<tbody>
	  <tr>
		<td><t t-esc="o.name"/></td>
		<td><t t-esc="o.age"/></td>
		<td><t t-esc="o.classroom_id.display_name"/></td>
	  </tr>
	</tbody>
  </table>
</div>
\end{lstlisting}
\end{frame}

\begin{frame}{QWeb Best Practices}
	\begin{itemize}
		\item Keep templates readable and modular with \texttt{t-call}.
		\item Prefer \texttt{t-esc}/\texttt{t-out} over \texttt{t-raw}.
		\item Put complex calculations in Python methods/fields.
		\item Reuse company layout templates for consistent branding.
		\item Test with records that have empty related fields.
	\end{itemize}
\end{frame}
